% !TEX root = ../main.tex
% ---------------------------------------------------------------------------
% Section 5: the unified model (conversational delayed streams).
% Built: the streaming ASR backbone (Sec. 5.1). Planned: Secs. 5.2-5.4; every
% planned part carries a \draftnote.
% Budget: <= 620 words of prose (wordcount.py) + Fig. 3, about 1.1 pages.
% Moved to App. F (app:model): encoder and thinker specifications, the loss
% weights, delay medians, backbone WERs, the [56,3] trade-off, Stage-0/1 and
% Stage-2/3 data recipes, per-speaker forced alignment before mixing, the
% decoding procedure and cap, the session length in audio positions,
% licences, EXAONE exclusion, the real-time context (also in Sec. 7).
% Moved to App. G (app:units): minimum-unit sizes, pilot LAAL/chrF values,
% collapse counts, the TAXI turn-length comparison.
% Stated once elsewhere, not repeated here: that the baselines, unlike the
% model and the M1-based controls, never see simulated sessions (Sec. 6.6,
% App. F); the decoupled twin (Sec. 6.6, App. D, App. F); that the output
% language fixes the direction (Def. 1). Coupling claims are left to the
% decoupled twin and the same-backbone cascade (Sec. 6.6).
% Within a turn only verified units are committed; a turn end commits the
% rest, so the novelty sentence says "commit units verified ...", not "only".
% Floats: Fig. 3 (figure*, fig:model), defined in sections/04_bench.tex so
% that it lands on the page where this section starts. One sentence per line.
% Labels kept: sec:model, sec:model-backbone, sec:model-stream,
% sec:model-targets and sec:model-inference (both on Sec. 5.3, which now
% holds the latency accounting), sec:model-training, eq:chunk, eq:loss.
% eq:delay was only referenced here; the delay rule is now inline.
% ---------------------------------------------------------------------------
% Special tokens of the serialized stream, e.g. \mdltok{next}, \mdltok{A}, \mdltok{{\to}de}.
% Robust, so it also works in captions; later files may reuse it.
\DeclareRobustCommand{\mdltok}[1]{\ensuremath{\langle\mathrm{#1}\rangle}}

\section{\model: Conversational Delayed Streams}\label{sec:model}

% Fig. 3 (figure*, fig:model) is defined in sections/04_bench.tex, before
% \paragraph{Scoring.}: a figure* appears at the earliest on the page after its
% definition, so defining it here would push it one page past this section.

After each $\Delta=80$\,ms audio chunk, the causal speech LLM \model writes zero or more tokens of a speaker-tagged transcript and translations into both languages (\cref{fig:model}): reading and writing become next-token prediction.
Unlike delayed streams modeling (one delayed text stream; \citealp{zeghidour2025streaming}), \emph{conversational delayed streams} delay each channel separately; unlike Hibiki's alignment-based per-word delays \citep{labiausse2025highfidelity}, they commit units verified by translating the source heard so far.
\draftnote{planned: only the backbone of \cref{sec:model-backbone} is built; \cref{sec:model-stream,sec:model-targets,sec:model-training} describe the design, and no \task result exists yet}

\subsection{Streaming backbone}\label{sec:model-backbone}
A cache-aware FastConformer encoder from Nemotron~3.5 ASR Streaming 0.6B \citep{nvidia2026nemotron35asr}, with a measured look-ahead of at most 80\,ms, feeds one embedding $a_k$ per chunk $k$, via an MLP adapter, to the 1.7B Qwen3-ASR language model, the \emph{thinker} \citep{shi2026qwen3asr}, replacing its non-causal audio encoder (\cref{app:model}).
After a prefix with the language and a delay token \mdltok{DELAY\_\delta}, each chunk holds $a_k$, the text tokens due in it and an end-of-chunk token \mdltok{next}; a word ending at $t$ is due in chunk $\lfloor t/\Delta\rfloor+\delta$.

\subsection{One serialized stream}\label{sec:model-stream}
The prefix names the language pair and two delays, $\delta_s$ for transcripts and $\delta_t$ for translations, sampled per training session.
Chunk $k$ reads
\begin{equation}\label{eq:chunk}
  B_k = a_k\,\mdltok{\hat s_1}\,z_1\cdots\mdltok{{\to}\lambda_1}\,y_1\cdots\mdltok{next},
\end{equation}
where $z_j$ are words of speaker $\hat s_j\in\{\mathrm{A},\mathrm{B}\}$, numbered by arrival, and $y_i$ translates one unit into language $\lambda_i$.
Each tagged segment is a committed piece (\cref{def:cst}), making the speaker and direction state explicit; under overlap, words are ordered by end time, as in t-SOT \citep{kanda2022streaming} and JSTAR \citep{moritz2025transcribing}.
We minimise
\begin{equation}\label{eq:loss}
  \mathcal{L}=-\textstyle\sum_{m\in\mathcal{Y}}w_m\log p_\theta(o_m\mid o_{<m})\,\big/\sum_{m\in\mathcal{Y}}w_m
\end{equation}
over the labelled positions $\mathcal{Y}$ (text, tags, \mdltok{next}) of the tokens $o$, with $w_m<1$ only for \mdltok{next} (\cref{app:model}); a two-slot, arrival-order activity head \citep{park2025sortformer} adds a binary cross-entropy and serves DER.
The session-long KV cache lets both sides condition every decision.
\draftnote{planned: the two speaker tags are reserved in the vocabulary but untrained; translation segments, the translation delay token and the activity head do not exist yet}

\subsection{Commit targets and latency}\label{sec:model-targets}\label{sec:model-inference}
A teacher LLM (Qwen3.8-27B) translates each source turn $u_{1:I}$ into $v$, and contextual alignment \citep{labiausse2025highfidelity} maps target word $v_j$ to $\alpha_j=\arg\max_i\,[\phi_j(i)-\phi_j(i-1)]$, $\phi_j(i)$ being its log-probability when $v$ is force-decoded from $u_{1:i}$.
Position $i$ is a candidate if $\{j:\alpha_j\le i\}=\{1,\dots,\kappa\}$ with $\kappa$ above the committed length, as meaning units require \citep{zhang2020learning}, and the source has grown by a minimum unit.
A candidate becomes a verified meaning unit (MU2) if the teacher's translation of the new source, given the committed text, is in the target language and reaches chrF${}\ge50$ against the newly aligned span of $v$.
Its translation is due $\delta_t\ge\delta_s$ chunks after its source's last chunk; a turn end commits the rest of the turn, so \model must detect it.
Relative to the teacher's alignment, targets thus never anticipate unheard speech, whereas sentence-level targets \citep{ma2019stacl} require anticipating clause-final German and Korean verbs; but commit points depend on the whole turn, so \model imitates a privileged teacher.
In a text-level \pair{Ko}{En} pilot (silver references; utterances much longer than \taxi turns), MU2 roughly halves the lag of utterance-end commits (a LAAL variant), losing at most 2.2 chrF (\cref{app:units}).
\draftnote{planned: MU2 targets for \pair{En}{De} and at training scale; the German minimum unit and the $\delta_t$ values are set on held-out data}

A token written after chunk $k$ is time-stamped $(k+1)\Delta+\ell_{\mathrm{enc}}+c$, so on schedule its latency after evidence ending at $t$ is $q+\delta\Delta+\ell_{\mathrm{enc}}+c$: quantisation $q\in(0,\Delta]$, stream delay $\delta\in\{\delta_s,\delta_t\}$, encoder look-ahead $\ell_{\mathrm{enc}}$ (80\,ms) and compute time $c$ (zero on the ideal clock; \cref{tab:ca-latency}).
As turn ends force commits, \eo then stays near $q+\delta_t\Delta+\ell_{\mathrm{enc}}+c$.

\subsection{Training without real cross-lingual conversations}\label{sec:model-training}
After single-speaker English and Korean ASR training (Stages~0--1), Stage~2 trains M1 for speaker-attributed transcription on simulated sessions with gaps and overlaps drawn more broadly than the benchmark's; Stage~3 adds translations committed at turn ends (M2) or at MU2 (M3) (\cref{tab:training}).
Training excludes \taxi and \syn dialogues and voices (\cref{app:model}).
\draftnote{planned: Stages 2--3; German training data has yet to be obtained}
